\section{Comparing local curve classes}\label{sec:comparison}

\subsection{Two normal weights}
A degree-$m$ cycle in the total space of $L_+\oplus L_-$, with
$L_+=L$ and $L_-=K_C L^{-1}$, meets a fiber of the projection to $C$
in a multiset of length $m$. No fiber contains a compact curve.
The elementary symmetric functions of either coordinate in this multiset
are sections of $L_\pm^j$, $1\le j\le m$. For instance, on the
normalization of a component the characteristic polynomial of
multiplication by the fiber coordinate on its finite flat direct image
extends the symmetric functions across branch values. Multiply the
polynomials with their cycle multiplicities. Evaluation at more than
$\max(0,m(g-1))$ distinct points therefore gives a tuple with common
zero precisely $m[C]$. Its weights are the integers $\pm j$.

For clarity, a tuple of $k$ nonzero weights $r_1,\ldots,r_k$ on the
cycle space gives, over the finite Borel base $\PP^k$, a section of the
pullback of $\bigoplus_i\cO(r_i)$. Pull back its relative Thom class and
divide by $\prod_i r_i$. The support on the relative cycle space is
compact. Adding an equivariant test after passage to the next projective
space gives the hyperplane pushforward of this class, since the normalized
Euler class of the added line is the hyperplane class. Concatenating two
tuples proves independence of the tests. Varying the curve and line of
degree $g-1$ gives the same universal class $a_{g,m}$
\cite[\S\S2--3]{WueBalanced}. The division in this prescription explains why
integrality of its answer requires proof.

The first two weights can already distinguish the prescriptions. Over a
curve of genus two the scalar and balanced doubles satisfy
$a_{2,2}-x_{2,2}=-2x_{4,1}$. In the holomorphic double-cover
calculation two smooth branches of genus four move away from the
zero section. Their contribution fixes the sign in the comparison
\cite[\S8]{WueProduct}.

There are two operations on the resulting coefficients. Pullback by
addition of effective cycles is the coproduct $\Delta$. Pullback by
multiplication of a cycle by $l$ is denoted $\rho_l$. With $u_0=1$, the
scalar and the antidiagonal classes of a fixed genus both satisfy
\begin{equation}\label{eq:localops}
 \Delta u_d=\sum_{i+j=d}u_i\otimes u_j,\qquad
 \rho_l u_d=\begin{cases}u_{d/l}&l\mid d,\\0&l\nmid d.\end{cases}
\end{equation}
These identities include the diagonal of Borel bases in the coproduct.
In each shifted term the two relative cohomological degrees are $d$
and $-d$, with $d\ne0$. In Chern number zero the factor of positive
relative cohomological degree has negative virtual dimension and vanishes
\cite[Lemma~4.13 and Theorem~6.13]{Pardon},
\cite[Proposition~3.1]{WueBalanced}. The augmentation is restriction to the
empty cycle.

The laws \eqref{eq:localops} alone leave a choice of integral
coordinates. For instance one may replace $A_g(q)$ by
$A_g(q)A_h(q^2)$ for $h\ne g$ and invert the change integrally.
The unit of a positive elementary curve specifies the particular
coordinates used in Theorem~\ref{thm:compact}.

\subsection{Elementary curves and multiplicity one}
We use the comparison $\Phi:H\to S$ of \cite[\S2]{WueProduct}, where $S$
is the rational cycle algebra for smooth families of almost-complex
six-manifolds. Over a point the characteristic function of a compact
open set of cycles gives an insertion. Open extension in the target,
external product for disjoint targets, and the two endpoint restrictions
of a compactly supported class over an interval are the operations we
will need. The geometric topology records the support and the integration
current of the cycle. In integrable families it gives the topology of
compact analytic cycles.

We recall three results from \cite{WueProduct} and their roles in the
comparison. First, on the elementary genus-$h$ curve of Ionel and Parker
the unit at its $j$-fold cycle is $\Phi(a_{h,j})$
\cite[Proposition~3.1]{WueProduct}. On the sum of two copies of a spin line,
real rotations commute with the Lee deformation to a super-rigid normal
operator; in complex eigenlines those rotations have weights $+1,-1$.
The normalized Thom class extends on the finite Borel family, restricting
at the elementary end to the unit of the isolated cycle. Second,
a bounded-energy partition into clusters gives a \emph{product} of these
units and their inverses \cite[\S4, Proposition~4.1]{WueProduct}. It is
essential here to partition disconnected cycles as well as connected
maps. We explain its local use in \S\ref{sec:basis}.

Third, the classes $\Phi(x_{h,1})$ are linearly independent over $\QQ$
\cite[Theorem~5.1]{WueProduct}. The argument tests a proposed finite
relation at its least genus $h$. The image-cycle map on connected stable
maps of that genus is proper over compact sets of cycles. A class supported
on an embedded curve of larger genus has no preimage in multiplicity one.
At genus $h$ the remaining regular embedded curve has a single preimage
with trivial automorphism group and complex local degree one. Stability
excludes trees of contracted rational components there. The
compact-support argument on a smooth parameter family representing a
putative relation promotes this local orientation calculation to
independence \cite[Lemmas~5.2--5.3 and Proposition~6.1]{WueProduct}.
Thus we need injectivity of $\Phi$ only on
$V=\bigoplus_h\QQ x_{h,1}$. The argument uses a fixed genus for each
test of a relation.

\begin{lemma}\label{lem:detect}
In $H$, the primitive elements annihilated by every $\rho_l$, $l>1$,
are precisely $V$. Two series with constant coefficient one satisfying
\eqref{eq:localops} through $q^m$ agree to that order if their
images under $\Phi$ do.
\end{lemma}
\begin{proof}
For the first assertion we recall the algebraic argument of
\cite[Lemma~7.1]{WueProduct}. The scalar classes generate $H$ rationally
by \cite[Theorem~6.13]{Pardon}. In the polynomial algebra on formal
scalar generators change variables to the coefficients $v_{h,j}$ of
the logarithm of $1+\sum_{j\ge1}x_{h,j}q^j$. These are primitive, and
$\rho_l v_{h,j}$ is $v_{h,j/l}$ if $l$ divides $j$, and zero otherwise.
These formulas already hold formally. A polynomial relation of smallest
total degree in the new variables has zero constant term. Its reduced
coproduct lies in the tensor square of the space of polynomials of
strictly smaller degree. Evaluation in $H$ is injective on that space,
hence also on its tensor square over $\QQ$. This makes the relation a
primitive polynomial, which is linear in characteristic zero by the
identity $f(v+w)=f(v)+f(w)$. Applying division by its largest label
$j$ would give a nonzero relation among the $x_{h,1}$ (use the identity
if $j=1$). Their smooth images are independent. This proves polynomial
independence of the scalar generators. The same linearity and largest-label
argument in $H$ now identifies its joint primitive kernel with $V$.
For the second assertion, the first coefficient of a difference is
primitive by addition and killed by every proper division. Its smooth
image vanishes, and $\Phi$ is injective on $V$.
\end{proof}
